%% Instructivo para estudiantes — Sistemas de Coordenadas (CB104)
%% Versión LaTeX, paso a paso, del INSTRUCTIVO_DE_USO.md de la raíz, recortada a lo que
%% necesita quien cursa la materia (sin la parte de adopción por otra cátedra ni la de
%% compilación del apunte). Usa el mismo formato que el apunte (apunte/apunte-formato.sty).
%% Compilar: instructivo/compilar.sh (lualatex vía latexmk, con TEXINPUTS apuntando a apunte/).
\documentclass[11pt,a4paper]{report}
\usepackage{apunte-formato}

%% Numeración de secciones sin el prefijo de capítulo (este documento no usa \chapter).
\renewcommand{\thesection}{\arabic{section}}
\renewcommand{\thesubsection}{\thesection.\arabic{subsection}}
\setcounter{secnumdepth}{2}
\setcounter{tocdepth}{1}

\renewcommand{\apunteEncabezadoIzq}{Sistemas de Coordenadas (CB104)}
\renewcommand{\apunteEncabezadoDer}{Instructivo para estudiantes}
\hypersetup{
  pdftitle={Sistemas de Coordenadas: instructivo para estudiantes},
  pdfauthor={Gustavo Varela},
  pdfsubject={Cómo usar el sitio, el apunte y el repositorio de Sistemas de Coordenadas (CB104, FIUBA)}
}

%% Comandos de terminal: mismo listings del apunte, pero sin resaltado de Python.
\lstdefinestyle{terminal}{language={},basicstyle=\small\ttfamily}
\newcommand{\clave}[1]{\texttt{\detokenize{#1}}}
%% Identificador de unidad en la tabla: enlace al sitio, monoespaciado chico.
\newcommand{\idUnidad}[1]{\href{\sitioURL/unidades/#1/index.html}{\footnotesize\ttfamily\detokenize{#1}}}

\begin{document}
\pagenumbering{arabic}\apunteCuerpotrue

%% --- Portada breve ---
\begin{center}
  {\color{muted}\large Sistemas de Coordenadas (CB104) \textperiodcentered{} Agrimensura \textperiodcentered{} FIUBA}\par\medskip
  {\color{tealdark}\Huge\bfseries Instructivo para estudiantes}\par\medskip
  {\Large Cómo usar el sitio, el apunte y el repositorio, paso a paso}\par\bigskip
  {\color{muted} Material del trabajo final de la materia \textperiodcentered{} Gustavo Varela \textperiodcentered{} octubre de 2026}
\end{center}

\begin{idea}[Qué hay y cómo se usa]
El material tiene cuatro piezas, todas en el mismo sitio web: veinte \textbf{unidades interactivas} (simulaciones que corren en el navegador), un \textbf{apunte} en PDF con las definiciones, las demostraciones y las guías de problemas resueltas, un \textbf{mapa de conceptos} para repasar y una \textbf{autoevaluación} de 120 preguntas.

Hay dos maneras de usarlo. La primera, que alcanza para toda la cursada, es desde el navegador, sin instalar nada (secciones \ref{sec:accesos} a \ref{sec:autoevaluacion}). La segunda es bajar el repositorio a tu computadora para abrir las unidades sin internet o para experimentar con el código (sección \ref{sec:local}). La sección \ref{sec:problemas} junta los tropiezos más comunes y cómo salir de cada uno.
\end{idea}

\makeatletter
\section*{Contenido}\@starttoc{toc}
\makeatother

%% ===================================================================
\section{Los tres accesos}\label{sec:accesos}

Todo corre en el navegador. No hace falta instalar nada ni tener cuenta en ningún lado.

\begin{resultado}[Direcciones]
\begin{itemize}
  \item \textbf{Curso interactivo} (portada con las veinte unidades): \url{\urlsitio}
  \item \textbf{Apunte en PDF} (definiciones, demostraciones y guías resueltas): \url{\sitioURL/apunte.pdf}
  \item \textbf{Mapa de conceptos}: \url{\sitioURL/mapa.html}
  \item \textbf{Autoevaluación} (120 preguntas, tres modos): \url{\sitioURL/autoevaluacion.html}
  \item \textbf{Repositorio} con todas las fuentes: \url{https://github.com/gusvarela-fiuba/Sistemas-de-Coordenadas}
\end{itemize}
\end{resultado}

\begin{atencion}[La dirección distingue mayúsculas]
El tramo \clave{Sistemas-de-Coordenadas} va con las tres iniciales en mayúscula. Escrito en minúsculas, el servidor responde \enquote{404 Not Found}. Si llegaste a esa pantalla, revisá la dirección antes que cualquier otra cosa.
\end{atencion}

%% ===================================================================
\section{La portada: por dónde entrar}\label{sec:portada}

La portada es el hilo conductor del curso: las veinte unidades en orden, cada una con su pregunta de campo, su gancho de agrimensura y los enlaces a la lección, a la ficha del capítulo y a las secciones del apunte que la respaldan. Arriba de todo hay una fila de entradas. Las que importan para estudiar son cuatro:

\begin{enumerate}
  \item \textbf{Empezar.} Lleva a la Unidad 0, \emph{El idioma de la materia}. Es la entrada si recién arrancás la cursada o si querés ponerte el vocabulario antes de seguir. Está pensada para el nivel de álgebra con el que se llega del CBC.
  \item \textbf{Seguir el programa.} Lleva a la unidad que corresponde al avance del curso, con la indicación de a cuál apunta. Es la entrada para ir al día con las clases.
  \item \textbf{Repasar.} Abre el mapa de conceptos completo. Tiene un selector «Ruta» («todo» o «solo troncal») y un deslizador «hasta el capítulo N» que atenúa los conceptos y las flechas posteriores al capítulo elegido: así ves qué sabés hasta hoy y qué viene después. Debajo del mapa está la lista de capítulos con sus etiquetas y enlaces.
  \item \textbf{Autoevaluación.} Abre la página de las 120 preguntas, con sus tres modos (sección \ref{sec:autoevaluacion}).
\end{enumerate}

Las otras entradas («Para el docente» y las dos «Para el aula») describen el material desde el lado de quien da la clase. Podés mirarlas, pero no las necesitás para estudiar.

Más abajo en la portada está el \textbf{Recorrido paso a paso} con las veinte unidades. Cada una se abre desde ahí.

%% ===================================================================
\section{Cómo usar una unidad, paso a paso}\label{sec:unidad}

Cada unidad se abre en una dirección de la forma \expandafter\nolinkurl\expandafter{\sitioURL/unidades/<id>/index.html}, donde \clave{<id>} es el identificador que figura en la tabla de la sección \ref{sec:recorrido}. Las lecciones corren adentro del navegador con WebAssembly: es Python ejecutándose en la página, sin servidor. La primera carga tarda unos segundos mientras baja el intérprete y las bibliotecas; después todo responde al instante.

Todas las unidades siguen el mismo molde de seis secciones, siempre de lo geométrico al cálculo. Conviene recorrerlas en este orden:

\begin{enumerate}[start=0]
  \item \textbf{Situación de campo y predicción.} Un problema concreto de agrimensura y una pregunta de opción múltiple que se responde sin hacer cuentas. Elegí una opción en serio antes de mirar la respuesta: cada opción incorrecta explica el error concreto que la produce, y eso es lo que más enseña.
  \item \textbf{Hecho geométrico.} Qué se conserva y qué cambia, antes de cualquier fórmula, con la escalera \(\R \to \R^2 \to \R^3\) en una tabla. Leela despacio: es la parte que da sentido a todo lo que sigue.
  \item \textbf{Manipulación.} La figura interactiva de tres capas: la grilla de antes (gris), la grilla después de la transformación (verde azulado) y un objeto reconocible (una parcela, un cubo). Arriba del deslizador hay casos precargados (identidad, 90°, un ángulo cualquiera y los contraejemplos: reflexión, corte, determinante cero). Al lado están los \enquote{números vivos} (área con signo, determinante) y la matriz actual, con cada columna pintada del color de la flecha que representa en la figura. El botón de reproducir anima el camino desde la identidad hasta la transformación elegida. Mové los controles hasta que puedas anticipar qué va a pasar antes de soltarlos.
  \item \textbf{Cuenta a mano y código lado a lado.} La misma cuenta en lápiz (con las columnas coloreadas) y en NumPy, una al lado de la otra. El código que ves es el mismo que corre en la celda.
  \item \textbf{Python de hoy.} Una sola idea nueva de Python por unidad, aplicada al problema de campo, con la forma de los arreglos y un ejemplo chico.
  \item \textbf{Ejercicios.} Dos de lápiz, con la respuesta escondida en un desplegable «Ver respuesta», y dos de código con hueco: completás los \clave{...} en un editor en vivo y una celda de abajo verifica el resultado contra la referencia. Te dice si coincide, si la forma del arreglo está mal o si todavía no corre. Cada ejercicio lleva un nivel (elemental, medio o avanzado).
\end{enumerate}

Al cierre, el anexo \textbf{\enquote{Dónde está la demostración}} enlaza las secciones del apunte que respaldan la unidad, más la ficha del capítulo, el hilo conductor y las unidades anterior y siguiente. La regla es simple: lo interactivo muestra, las demostraciones están en el apunte.

\begin{atencion}[Qué se puede tocar y qué no]
En las lecciones publicadas no se pueden editar las celdas del cuaderno: corren en modo aplicación. Lo editable son los controles (deslizadores, menús, botones) y los editores de los ejercicios de código. Para modificar una lección entera hay que bajar el repositorio (sección \ref{sec:local}).
\end{atencion}

%% ===================================================================
\section{El recorrido de las veinte unidades}\label{sec:recorrido}

El orden es el del hilo conductor de la portada. La columna \enquote{Apunte} dice qué capítulos del PDF respaldan cada unidad; la columna \clave{<id>} es la que necesitás para armar la dirección de una unidad o para abrirla en tu computadora. Ojo: el número de unidad y el prefijo del identificador no coinciden.

\begingroup\small
\begin{longtable}{@{}p{13mm}p{50mm}p{58mm}p{32mm}@{}}
\toprule
\textbf{Unidad} & \textbf{Título (enlace)} & \clave{<id>} & \textbf{Apunte}\\
\midrule
\endfirsthead
\toprule
\textbf{Unidad} & \textbf{Título (enlace)} & \clave{<id>} & \textbf{Apunte}\\
\midrule
\endhead
\bottomrule
\endfoot
0 & \linkUnidad{00_fundamentos}{El idioma de la materia} & \idUnidad{00_fundamentos} & caps.\ 1, 5 y 10 (complementario)\\
1 & \linkUnidad{01a_punto_quieto}{Rotaciones I: el punto no se movió} & \idUnidad{01a_punto_quieto} & cap.\ 1\\
2 & \linkUnidad{01b_rotacion_plano}{Rotaciones II: la rotación en el plano} & \idUnidad{01b_rotacion_plano} & cap.\ 1\\
3 & \linkUnidad{01c_rotaciones_3d}{Rotaciones III: tres ejes y el orden importa} & \idUnidad{01c_rotaciones_3d} & caps.\ 1 y 2\\
4 & \linkUnidad{02_producto_interno}{Producto interno: proyección, longitud, ángulo} & \idUnidad{02_producto_interno} & cap.\ 1\\
5 & \linkUnidad{03_reflexiones_similaridades}{Reflexiones y similaridades} & \idUnidad{03_reflexiones_similaridades} & caps.\ 2 a 4\\
6 & \linkUnidad{04_afines_homogeneas}{Afinidades y coordenadas homogéneas} & \idUnidad{04_afines_homogeneas} & caps.\ 3 y 4\\
7 & \linkUnidad{05_proyectivas}{Proyectivas: la foto del dron y el pentágono} & \idUnidad{05_proyectivas} & complemento C1\\
8 & \linkUnidad{06_formas_lineales}{Formas lineales y dualidad} & \idUnidad{06_formas_lineales} & cap.\ 5\\
9 & \linkUnidad{07a_cuna_determinante}{Cuña y determinante: el detector de colinealidad} & \idUnidad{07a_cuna_determinante} & cap.\ 6\\
10 & \linkUnidad{07b_nucleo_autovectores}{Núcleo y autovectores: las direcciones que no giran} & \idUnidad{07b_nucleo_autovectores} & caps.\ 6 y 7\\
11 & \linkUnidad{07c_autovalores_3d_puntos_fijos}{Autovalores en 3D y puntos fijos proyectivos} & \idUnidad{07c_autovalores_3d_puntos_fijos} & cap.\ 7\\
12 & \linkUnidad{08_cambio_de_base}{Cambio de base y diagonalización} & \idUnidad{08_cambio_de_base} & cap.\ 9\\
13 & \linkUnidad{09_cambio_de_base_svd}{Cambio de base, SVD e inversas} & \idUnidad{09_cambio_de_base_svd} & caps.\ 8 y 11\\
14 & \linkUnidad{10_jacobiano_primera_forma}{Jacobiano y primera forma fundamental} & \idUnidad{10_jacobiano_primera_forma} & cap.\ 10\\
15 & \linkUnidad{11_cuatro_problemas}{Cuatro problemas: regresión, círculos, Procrustes, regresión total} & \idUnidad{11_cuatro_problemas} & cap.\ 12\\
16 & \linkUnidad{12_curvatura}{Curvatura como forma cuadrática} & \idUnidad{12_curvatura} & cap.\ 13\\
17 & \linkUnidad{13_elipsoide}{Geometría del elipsoide} & \idUnidad{13_elipsoide} & cap.\ 14\\
18 & \linkUnidad{14_trigonometria_esferica}{Trigonometría esférica} & \idUnidad{14_trigonometria_esferica} & cap.\ 15\\
19 & \linkUnidad{15_transformaciones_conformes}{Transformaciones conformes: del elipsoide al plano cartográfico} & \idUnidad{15_transformaciones_conformes} & caps.\ 14 y 15 (complementario)\\
\end{longtable}
\endgroup

\subsection{Cómo se encadenan}

Cada unidad declara de qué ladrillo viene y hacia cuál va; la propia lección lo muestra al principio como «venimos de / vamos hacia». El edificio, resumido en cinco tramos:

\begin{itemize}
  \item \textbf{Unidades 0 a 3.} Con el idioma puesto (grupos, escalares, punto y vector, formas lineales, métrica), la primera transformación: la rotación. Primero giran los ejes y el punto se queda quieto; después la rotación como matriz y como número complejo; después los tres ejes de \(\R^3\), donde el orden importa.
  \item \textbf{Unidades 4 a 7.} Lo que la rotación conserva se mide con el producto interno. Con eso se clasifican las reflexiones y las similaridades, se meten las traslaciones adentro de la matriz con coordenadas homogéneas y, al liberar la última fila, aparecen las proyectivas: las paralelas se cruzan como en una foto.
  \item \textbf{Unidades 8 a 11.} Las rectas se transforman al revés que los puntos; para entenderlo hace falta el dual. De ahí salen el determinante como área orientada, el núcleo cuando el determinante se anula, los autovectores como direcciones que no giran y, en \(\R^3\), el eje de una rotación.
  \item \textbf{Unidades 12 a 15.} Con los autovectores como ejes nuevos la matriz se vuelve diagonal. Cuando esa grilla propia no existe, dos grillas ortogonales siempre alcanzan: la SVD. El terreno es una transformación que cambia de punto a punto (el jacobiano), y la misma tabla de productos internos reaparece en las ecuaciones normales de los cuatro problemas de ajuste.
  \item \textbf{Unidades 16 a 19.} Cuánto se aparta una superficie de su plano tangente (curvatura), una superficie concreta (el elipsoide), triángulos sobre ella (trigonometría esférica) y cómo volcarla al papel sin torcer los ángulos (transformaciones conformes). El arco se cierra en los números complejos de la Unidad 2.
\end{itemize}

%% ===================================================================
\section{El apunte en PDF}\label{sec:apunte}

El apunte es el respaldo formal de las unidades: definiciones, teoremas con demostración y guías de problemas resueltas. Se baja o se abre en \url{\sitioURL/apunte.pdf}.

\begin{itemize}
  \item \textbf{Cómo está armado.} Un capítulo 0 con el idioma de la materia y la tabla de notación; los capítulos 1 a 15 siguen el programa; cuatro complementos (C1 a C4) con temas que se apoyan en varios capítulos; y los apéndices, con las respuestas de los problemas graduados que cierran cada capítulo.
  \item \textbf{Cómo leer los enlaces.} El símbolo \(\triangleright\) marca una unidad interactiva y \(\circlearrowright\) una animación; los dos son hipervínculos al sitio. El texto en \texttt{monoespaciado} es un identificador de archivo o una ruta del sitio.
  \item \textbf{Cómo combinarlo con las unidades.} Un camino que funciona: abrir la unidad, hacer la predicción y la manipulación, y recién después leer en el apunte la sección que el anexo \enquote{Dónde está la demostración} señala. La cuenta de la unidad y la demostración del apunte usan la misma notación.
  \item \textbf{Las seis animaciones.} Hay seis animaciones cortas en bucle, enlazadas desde el apunte y desde el sitio, para mirar un concepto moverse (por ejemplo, la composición de rotaciones que no conmuta). No reemplazan la manipulación de la unidad: muestran un recorrido fijo.
\end{itemize}

%% ===================================================================
\section{La autoevaluación}\label{sec:autoevaluacion}

Son 120 preguntas de opción múltiple: seis por unidad, dos elementales, dos medias y dos avanzadas. Cada pregunta tiene cuatro opciones, una sola correcta, y una explicación para cada opción. Aparecen al final de cada unidad y, todas juntas, en \url{\sitioURL/autoevaluacion.html}, con tres modos:

\begin{enumerate}
  \item \textbf{Práctica.} Elegís una unidad (o todas) y un nivel (o todos). Al marcar una opción ves al instante si acertaste y por qué; si fallaste, también cuál era la correcta. El botón «Mezclar» vuelve a barajar preguntas y opciones. No guarda nada. Es el modo para estudiar.
  \item \textbf{Simulacro.} Veinte preguntas, una por unidad y en el orden de las unidades, con el nivel sorteado. No hay devolución hasta que apretás «Entregar»; entonces ves el puntaje total, el puntaje por nivel y la corrección completa con explicaciones. «Nuevo simulacro» sortea otro. No guarda nada. Es el modo para medirte antes de un parcial.
  \item \textbf{Evaluación (4 intentos).} Treinta preguntas: una por unidad más diez al azar. Al entregar ves el puntaje y qué unidades acertaste, sin explicaciones. Cada intento queda anotado (fecha, puntaje y nivel de cada pregunta) y podés hacer cuatro. Al agotarlos queda el historial, la mejor nota y un botón para borrarlo, que pide un segundo clic para confirmar. También podés descargar el resultado como archivo de texto.
\end{enumerate}

\begin{atencion}[El registro vive en tu navegador]
El historial de la evaluación se guarda solo en el navegador donde la hiciste. No se envía a ningún servidor, no lo ve nadie más, y se pierde si borrás los datos del sitio o cambiás de navegador o de computadora. Si querés dejar constancia, usá «Descargar resultado» al terminar cada intento.
\end{atencion}

%% ===================================================================
\section{Usar el repositorio en tu computadora}\label{sec:local}

Esta parte es opcional. Sirve para abrir las unidades sin internet, para que respondan más rápido que en el navegador, o para experimentar con el código de una lección. Funciona en Windows, Mac y Linux. Hacen falta dos programas: Git, para bajar el repositorio, y \clave{uv}, que instala Python y las bibliotecas por su cuenta sin tocar el Python de tu sistema.

\subsection{Paso 1: instalar Git y uv}

\begin{itemize}
  \item \textbf{Git}: \url{https://git-scm.com/downloads}. En Mac y Linux suele venir instalado o se instala con el gestor de paquetes; en Windows, con el instalador de esa página.
  \item \textbf{uv}: \url{https://docs.astral.sh/uv/getting-started/installation/}. En Mac o Linux, en una terminal:
\begin{lstlisting}[style=terminal]
curl -LsSf https://astral.sh/uv/install.sh | sh
\end{lstlisting}
  En Windows, en PowerShell:
\begin{lstlisting}[style=terminal]
powershell -ExecutionPolicy ByPass -c "irm https://astral.sh/uv/install.ps1 | iex"
\end{lstlisting}
  Después de instalar, cerrá la terminal y abrila de nuevo, para que el comando \clave{uv} quede disponible.
\end{itemize}

Para comprobar que quedó todo, en una terminal nueva:
\begin{lstlisting}[style=terminal]
git --version
uv --version
\end{lstlisting}
Los dos tienen que responder con un número de versión.

\subsection{Paso 2: bajar el repositorio}

En la carpeta donde quieras tenerlo:
\begin{lstlisting}[style=terminal]
git clone https://github.com/gusvarela-fiuba/Sistemas-de-Coordenadas.git
cd Sistemas-de-Coordenadas
\end{lstlisting}
Esto crea la carpeta \file{Sistemas-de-Coordenadas} con todo el material: las unidades en \file{unidades/}, el apunte ya compilado en \file{apunte/Apunte_Sistemas_de_Coordenadas.pdf}, las animaciones en \file{animaciones/} y las resoluciones de ejercicios en \file{ejercicios/}.

\subsection{Paso 3: instalar las dependencias}

Dentro de la carpeta:
\begin{lstlisting}[style=terminal]
uv sync
\end{lstlisting}
Crea un entorno aislado adentro del repositorio (la carpeta oculta \file{.venv}) con la versión de Python y las bibliotecas exactas que el material necesita. No instala nada fuera de la carpeta. La primera vez puede tardar un par de minutos porque baja Python y las bibliotecas; las siguientes veces es inmediato.

\subsection{Paso 4: abrir una unidad}

Cada lección es un archivo \file{unidades/<id>.py}, con el \clave{<id>} de la tabla de la sección \ref{sec:recorrido}. Para usarla tal como está en el sitio:
\begin{lstlisting}[style=terminal]
uv run marimo run unidades/08_cambio_de_base.py
\end{lstlisting}
Se abre el navegador con la unidad corriendo sobre el Python de tu computadora, sin mostrar las celdas de código y sin modificar el archivo. Para cerrarla, volvé a la terminal y apretá \clave{Ctrl+C}.

\subsection{Paso 5: experimentar con una copia}

Si querés ver el código de una lección y cambiarlo (los números de la situación de campo, los casos precargados, los ejercicios), trabajá sobre una copia cuyo nombre empiece con guion bajo:
\begin{lstlisting}[style=terminal]
cp unidades/08_cambio_de_base.py unidades/_ensayo_08.py
uv run marimo edit unidades/_ensayo_08.py
\end{lstlisting}
Se abre marimo con la lección completa, celdas incluidas, y cada cambio se ve al momento. La copia tiene que quedar dentro de \file{unidades/} para encontrar la biblioteca común \file{cb104.py}. El guion bajo inicial es la convención para reconocer a simple vista que es un borrador: el sitio nunca lo exporta y el original queda intacto.

\subsection{Paso 6: el botón de pánico}

Si tocaste un archivo original por error y la unidad dejó de andar:
\begin{lstlisting}[style=terminal]
git checkout -- unidades/08_cambio_de_base.py
\end{lstlisting}
Descarta tus cambios locales en ese archivo y restaura la versión del repositorio. Para volver todo el repositorio al estado publicado, \clave{git checkout -- .} desde la carpeta raíz. Las copias con guion bajo no se tocan con estos comandos; si querés sacarlas, borralas como cualquier archivo.

\subsection{Paso 7: levantar el sitio completo (optativo)}

Para tener en tu computadora la misma portada, el mapa y las veinte unidades que ves en línea:
\begin{lstlisting}[style=terminal]
uv run python tools/build.py
uv run python -m http.server -d site 8765
\end{lstlisting}
y abrí \url{http://localhost:8765/} en el navegador. El primer comando genera la portada, el mapa, las fichas, copia el apunte y exporta las veinte lecciones a WebAssembly en la carpeta \file{site/}; tarda entre uno y dos segundos por lección, y la primera vez después de \clave{uv sync} puede llegar a 40 segundos. El segundo sirve esa carpeta por HTTP, que es imprescindible: las lecciones exportadas no funcionan abriendo el archivo \file{index.html} directamente desde el explorador de archivos.

\subsection{Paso 8: mantenerlo al día}

El material se corrige y crece durante la cursada. Para traer la última versión:
\begin{lstlisting}[style=terminal]
git pull
uv sync
\end{lstlisting}
Si Git se queja de cambios locales, es porque editaste algún archivo original: aplicá antes el botón de pánico del paso 6 (tus copias con guion bajo no molestan) y repetí \clave{git pull}.

\begin{resultado}[Los ocho pasos en una pantalla]
\begin{lstlisting}[style=terminal]
git clone https://github.com/gusvarela-fiuba/Sistemas-de-Coordenadas.git
cd Sistemas-de-Coordenadas
uv sync                                    # una vez, y tras cada git pull
uv run marimo run unidades/<id>.py         # usar una unidad
cp unidades/<id>.py unidades/_ensayo.py    # copia para experimentar
uv run marimo edit unidades/_ensayo.py     # editar la copia
git checkout -- unidades/<id>.py           # deshacer cambios en un original
uv run python tools/build.py               # generar el sitio completo
uv run python -m http.server -d site 8765  # servirlo en localhost:8765
git pull && uv sync                        # actualizar
\end{lstlisting}
\end{resultado}

%% ===================================================================
\section{Problemas frecuentes}\label{sec:problemas}

\begin{tabularx}{\linewidth}{@{}>{\raggedright\arraybackslash}p{52mm}X@{}}
\toprule
\textbf{Qué pasa} & \textbf{Qué hacer}\\
\midrule
La dirección responde \enquote{404 Not Found}. & Revisá las mayúsculas de \clave{Sistemas-de-Coordenadas} (sección \ref{sec:accesos}).\\
Una unidad muestra la página en blanco o un indicador de carga durante varios segundos. & Es la primera carga del intérprete de Python en el navegador. Esperá; si pasa un minuto sin cambios, recargá la página. Las cargas siguientes son inmediatas.\\
Los controles no responden después de dejar la pestaña abierta mucho tiempo. & Recargá la página. No se pierde nada: las unidades no guardan estado.\\
El ejercicio de código dice que la forma del arreglo está mal. & El resultado tiene otra cantidad de filas o columnas que la referencia. Revisá si tenías que devolver un vector o una matriz, y si transpusiste de más.\\
El historial de la evaluación desapareció. & Vive en el navegador (sección \ref{sec:autoevaluacion}). Se pierde al borrar los datos del sitio o al cambiar de navegador o de computadora. Para la próxima, descargá el resultado al terminar.\\
\clave{uv: command not found} después de instalarlo. & Cerrá la terminal y abrí una nueva. Si persiste, volvé a correr el instalador de la sección \ref{sec:local}, paso 1.\\
La lección local abre pero no encuentra \clave{cb104}. & La copia de ensayo quedó fuera de \file{unidades/}. Movela adentro de esa carpeta.\\
El sitio generado no muestra las unidades al abrir \file{site/index.html}. & Hay que servirlo por HTTP con el segundo comando del paso 7. Abrirlo como archivo no funciona.\\
\clave{git pull} se niega por cambios locales. & Editaste un original. \clave{git checkout -- <archivo>} y repetí el \clave{git pull} (paso 8).\\
\bottomrule
\end{tabularx}

\medskip
Si encontrás un error en el material (una cuenta, un enunciado, un enlace roto), la forma de avisar es abrir un \emph{issue} en el repositorio: \url{https://github.com/gusvarela-fiuba/Sistemas-de-Coordenadas/issues}. Indicá la unidad o la página del apunte y qué esperabas ver.

\bigskip
{\small\color{muted}
El texto, el apunte, las figuras y las preguntas se publican con licencia CC BY-SA 4.0; el código, con licencia MIT. El detalle está en el archivo \file{LICENSE} del repositorio.\par}

\end{document}
